\documentclass[a4paper]{article}
\usepackage{amsfonts}
\usepackage{amsmath}
\usepackage{amssymb}
\usepackage{graphicx}     
\begin{document}
  
\noindent \large Auteur: Siebe Haché \\
\noindent \large Studierichting: Informatica\\
\noindent \large Oefeningenreeks 1C nr. 48.8\\
\medskip
\normalsize
$z^6 = -\sqrt{3} - i$\\

$r = \sqrt{(-\sqrt{3})^2 + (-1)^2} = \sqrt{3 + 1} = 2$\\

$\tan\theta = \dfrac{-1}{-\sqrt{3}} = \dfrac{1}{\sqrt{3}}$\\

we doen + $\pi$ omdat $-\sqrt{3} - i$ in het derde kwadrant zit\\

$\theta = \pi + \dfrac{\pi}{6} = \dfrac{7\pi}{6}$\\

$-\sqrt{3} - i = 2\cos\dfrac{7\pi}{6} + 2i\sin\dfrac{7\pi}{6}$\\

$w_k = \sqrt[6]{2}\left(\cos\dfrac{\dfrac{7\pi}{6} + 2k\pi}{6} + i\sin\dfrac{\dfrac{7\pi}{6} + 2k\pi}{6}\right)$\\

$w_0 = \sqrt[6]{2}\left(\cos\dfrac{7\pi}{36} + i\sin\dfrac{7\pi}{36}\right)$\\

$w_1 = \sqrt[6]{2}\left(\cos\dfrac{19\pi}{36} + i\sin\dfrac{19\pi}{36}\right)$\\

$w_2 = \sqrt[6]{2}\left(\cos\dfrac{31\pi}{36} + i\sin\dfrac{31\pi}{36}\right)$\\

$w_3 = \sqrt[6]{2}\left(\cos\dfrac{43\pi}{36} + i\sin\dfrac{43\pi}{36}\right)$\\

$w_4 = \sqrt[6]{2}\left(\cos\dfrac{55\pi}{36} + i\sin\dfrac{55\pi}{36}\right)$\\

$w_5 = \sqrt[6]{2}\left(\cos\dfrac{67\pi}{36} + i\sin\dfrac{67\pi}{36}\right)$\\
\begin{thebibliography}{999}
%\bibitem{a} Referentie
\end{thebibliography}
\end{document}